\section{Resultados de RT-2: ver la novelty y también los límites}

La parte de resultados presenta cuatro tipos de evidencia: casos de éxito, resúmenes por categoría, ablation y benchmarks específicos. Una lectura rigurosa pregunta por separado qué responde cada uno, en lugar de reunir todas las demos llamativas bajo la etiqueta de «emergence». Esta sección los desglosa en cuatro niveles: qualitative, quantitative, causal y compositional.

\subsection{Emergent skills: combinar instrucciones nuevas, no crear de la nada}

La qualitative grid muestra a RT-2 ejecutando acciones sobre objetos, símbolos y combinaciones semánticas que no aparecen directamente en los robot data, por ejemplo reconocer una marca de refresco de cola, elegir el objeto más pequeño o entender categorías abstractas. Lo que se llama emergent se parece más a \textbf{una nueva combinación de web semantics + learned motor primitive}; el espacio de acciones sigue proviniendo del robot controller entrenado y del token vocabulary.

\lecturefigure{slide-30-rt2-emergent-skills.jpg}{Emergent skills de RT-2: símbolos, sentido común y combinación de instrucciones nuevas}{Diapositiva V030 recuperada de la grabación oficial; 00:49:30}

\lecturefigure{slide-31-rt2-emergent-skills-chart.jpg}{Clasificación de las emergent skills: symbol understanding, reasoning, human recognition}{Diapositiva V031 recuperada de la grabación oficial; 00:50:24}

\begin{knowledgebox}{Lectura de la figura: qué oculta una gráfica de barras de promedios}
La gráfica de barras agrega varias task families; primero se comparan los modelos dentro de la misma category y después se mira el average. Si reasoning mejora mientras motor execution no cambia, el semantic prior está actuando; si human recognition es alto pero unseen environment es bajo, la transferencia de conceptos visuales no equivale a robustez ante la escena. El promedio no indica en qué objetos, acciones o entornos se concentran las fallas.
\end{knowledgebox}

\subsection{Clasificación seen / unseen y ablation}

La quantitative slide separa unseen objects, unseen backgrounds y unseen environments, y este diseño es necesario: el cambio de objeto pone a prueba la semantic recognition, el cambio de fondo pone a prueba la robustez visual y el cambio de entorno además modifica la geometry y la affordance. La ablation slide compara después co-fine-tuning, solo robot data y distintas model scale, para intentar identificar de dónde proviene la transfer.

\lecturefigure{slide-32-rt2-quantitative-generalization.jpg}{Tasa de éxito de RT-2 en el entorno seen y en varias configuraciones unseen}{Diapositiva V032 recuperada de la grabación oficial; 00:51:00}

\lecturefigure{slide-33-rt2-ablation-results.jpg}{Ablation de RT-2: contribución del co-fine-tuning y de la data mixture}{Diapositiva V033 recuperada de la grabación oficial; 00:51:48}

\begin{warningbox}{La ablation también tiene límites de extrapolación}
Si la ablation se realiza con el mismo robot, la misma camera, el mismo task template y el mismo evaluation protocol, puede mostrar el efecto de la configuración de entrenamiento sobre ese sistema, pero no responde directamente sobre hardware nuevo, control táctil, entornos humanos dinámicos ni tareas de varias horas. Cuanto más estricto es el control experimental, más clara es la explicación causal; al mismo tiempo, el ámbito de aplicación también debe declararse con claridad.
\end{warningbox}

\subsection{Language Table y chain-of-thought}

El Language Table benchmark evalúa manipulation sobre una mesa descrita mediante lenguaje, y sirve para comparar instruction following y generalización composicional. La página de chain-of-thought, por su parte, hace que el modelo genere primero un razonamiento semántico visual y después la action; muestra que el intermediate reasoning puede ayudar con tareas del tipo «mover a cierto lugar el objeto que puede contener agua», pero para saber si el reasoning textual impulsa realmente la acción de forma causal hace falta una intervention, no basta con mirar la trace en lenguaje natural.

\lecturefigure{slide-34-language-table-benchmark.jpg}{Language Table: resultados de generalización de las push-object instruction}{Diapositiva V034 recuperada de la grabación oficial; 00:52:03}

\lecturefigure{slide-35-rt2-chain-of-thought.jpg}{Chain-of-thought y secuencia de acciones de RT-2-PaLM-E}{Diapositiva V035 recuperada de la grabación oficial; 00:53:30}

\begin{knowledgebox}{Qué significa que «el reasoning ayudó a la policy»}
La evidencia más sólida no es que el modelo escriba una explicación aparentemente razonable, sino que modificar o retirar el intermediate reasoning cambie de forma sistemática el action success; o bien que el model supere claramente, en held-out tasks que requieren combinar relaciones en varios pasos, a un baseline de la misma capacidad que no genera reasoning. La trace en lenguaje natural es una ventana de diagnóstico; no equivale automáticamente a una faithful internal computation.
\end{knowledgebox}

\subsection{Síntesis de los VLA y problemas aún sin resolver}

La página de VLA summary enumera como beneficios las tareas nuevas, los objetos nuevos, la chain-of-thought y la generalization, y presenta como líneas futuras la data diversity, el context adicional y la ampliación de capability. No declara resuelto el low-level control; al contrario, los resultados indican que ya aparece la semantic transfer, pero los robot data, la precisión de las acciones y la closed-loop reliability siguen siendo los cuellos de botella decisivos.

\lecturefigure{slide-36-vla-summary.jpg}{Síntesis de capacidades de los modelos VLA y problemas futuros}{Diapositiva V036 recuperada de la grabación oficial; 00:55:36}

\teachervoice{El ponente llama a estos resultados «semantic transfer» y no inteligencia robótica general. El tono de la clase importa: que el modelo pueda ejecutar órdenes que no aparecieron textualmente en el entrenamiento es una capacidad nueva que merece atención; pero sigue operando en el mismo action space, la misma plataforma robótica y escenarios limitados.}

\subsection{Resumen de la sección}

Las conclusiones confiables de RT-2 se sostienen en varios niveles de evidencia a la vez: los qualitative cases muestran combinaciones nuevas, las category metrics miden la generalización, la ablation rastrea el co-fine-tuning y Language Table mide la instruction generalization. En conjunto respaldan que el VLM prior puede transferirse al control, pero no eliminan los límites de embodiment, confiabilidad y seguridad.
